\documentclass[../DoAn.tex]{subfiles}
\begin{document}

Chương trước đã trình bày phương pháp refinement pose được đề xuất nhằm cải thiện tính ổn định và tính hợp lý của chuyển động người 3D. Trên cơ sở đó, chương này tập trung mô tả quy trình thực nghiệm và phương pháp đánh giá hệ thống.

Nội dung chương bao gồm các chỉ số đánh giá được sử dụng trong bài toán human pose estimation, phương pháp xây dựng kịch bản thực nghiệm, các bộ dữ liệu sử dụng và cách đánh giá ảnh hưởng của từng thành phần refinement trong hệ thống. Phần này tập trung mô tả khung đánh giá để phục vụ việc thu thập và đối chiếu kết quả thực nghiệm ở giai đoạn sau.

\section{Các tham số đánh giá}

Để đánh giá hiệu quả của hệ thống refinement pose, đồ án sử dụng các chỉ số phổ biến trong lĩnh vực 3D Human Pose Estimation và Human Motion Analysis.

\subsection{MPJPE}

MPJPE (Mean Per Joint Position Error) là sai số khoảng cách Euclidean trung bình giữa vị trí khớp dự đoán và vị trí khớp ground truth trong không gian 3D.

\begin{equation}
\text{MPJPE} =
\frac{1}{N}
\sum_{i=1}^{N}
\left\|
P_i - \hat{P}_i
\right\|_2
\end{equation}

Trong đó:

\begin{itemize}
    \item $P_i$ là vị trí ground truth của khớp thứ $i$.
    \item $\hat{P}_i$ là vị trí khớp được dự đoán.
    \item $N$ là tổng số lượng khớp.
\end{itemize}

MPJPE được tính theo đơn vị milimét. Giá trị MPJPE càng nhỏ cho thấy sai số pose reconstruction càng thấp.

\subsection{PA-MPJPE}

PA-MPJPE (Procrustes Aligned MPJPE) là biến thể của chỉ số MPJPE sau khi thực hiện phép căn chỉnh rigid transformation giữa pose dự đoán và pose ground truth thông qua thuật toán Umeyama. Quá trình căn chỉnh bao gồm việc đồng nhất vị trí, hướng quay và tỷ lệ giữa hai pose thông qua các phép tịnh tiến (translation alignment), xoay (rotation alignment) và scale alignment. 

Nhờ loại bỏ các sai lệch toàn cục liên quan đến hệ tọa độ và kích thước cơ thể, PA-MPJPE tập trung đánh giá chất lượng tương đối của cấu trúc pose và khả năng tái tạo hình học của hệ thống.

\subsection{PCK}

PCK (Percentage of Correct Keypoints) đo tỷ lệ phần trăm các khớp được dự đoán nằm trong một ngưỡng khoảng cách xác định so với vị trí ground truth. Khoảng cách giữa khớp dự đoán và khớp ground truth được tính bằng khoảng cách Euclidean trong không gian 3D. Một khớp được xem là dự đoán chính xác nếu khoảng cách này nhỏ hơn ngưỡng đánh giá được thiết lập trước. 

Trong đồ án này, ngưỡng khoảng cách được lựa chọn theo các tiêu chuẩn phổ biến trong bài toán 3D pose estimation. Giá trị PCK càng cao cho thấy số lượng khớp được dự đoán chính xác càng lớn và chất lượng pose reconstruction càng tốt.

\begin{table}[H]
\centering
\caption{Các tiêu chí đánh giá sử dụng trong thực nghiệm}
\label{tab:evaluation_metrics_overview}
\begin{tabular}{|p{3.3cm}|p{4.6cm}|p{5.4cm}|}
\hline
\textbf{Chỉ số} & \textbf{Ý nghĩa} & \textbf{Dữ liệu cần có} \\
\hline
MPJPE / PA-MPJPE & Sai số hình học của tọa độ khớp, trước và sau Procrustes alignment & Ground truth pose 3D \\
\hline
PCK & Tỷ lệ khớp nằm trong ngưỡng sai số cho phép & Ground truth pose 3D \\
\hline
Temporal jitter & Mức rung giật hoặc biến đổi bất thường theo thời gian & Chuỗi pose theo thời gian \\
\hline
Bone-length variance & Mức ổn định cấu trúc xương trong chuỗi chuyển động & Skeleton và pose theo thời gian \\
\hline
Cross-view disagreement & Mức mâu thuẫn giữa hai view sau khi căn chỉnh & Hai chuỗi pose đã căn chỉnh \\
\hline
Runtime/frame & Khả năng triển khai và chi phí xử lý mỗi frame & Log thời gian xử lý \\
\hline
\end{tabular}
\end{table}

\section{Phương pháp thí nghiệm}

\subsection{Tập dữ liệu sử dụng}

Đồ án sử dụng hai nhóm dữ liệu chính nhằm đánh giá hệ thống trong các điều kiện chuyển động khác nhau.

MPI-INF-3DHP là bộ dữ liệu phổ biến trong bài toán monocular 3D pose estimation. Bộ dữ liệu này cung cấp ground truth 3D chính xác cùng nhiều góc quay và hành động khác nhau. MPI-INF-3DHP được sử dụng để đánh giá độ chính xác hình học của pose reconstruction trong các điều kiện chuyển động thông thường và ở môi trường phòng thí nghiệm.

Bên cạnh đó, đồ án sử dụng thêm bộ dữ liệu AthletePose3D nhằm đánh giá hệ thống trong các trường hợp chuyển động có tính động lực học cao. AthletePose3D bao gồm nhiều động tác thể thao phức tạp với tốc độ chuyển động lớn, thay đổi tư thế liên tục và xuất hiện nhiều trường hợp che khuất. Bộ dữ liệu này được sử dụng để phân tích khả năng duy trì tính ổn định của hệ thống refinement trong các điều kiện gần với môi trường thực tế hơn so với các bộ dữ liệu indoor thông thường.

\subsection{Pipeline thực nghiệm}

Pipeline thực nghiệm được xây dựng theo hướng postprocessing refinement cho bài toán tái tạo chuyển động người 3D từ video RGB. Hệ thống không huấn luyện lại mô hình pose estimation mà sử dụng trực tiếp kết quả đầu ra thô từ mô hình monocular pose estimation WHAM để thực hiện các bước refinement và fusion.

Đầu vào của hệ thống bao gồm hai video RGB được ghi từ hai góc nhìn khác nhau cùng với các dữ liệu đầu ra từ WHAM tương ứng với từng video. Các dữ liệu này bao gồm mesh vertices, tham số pose, tham số shape, translation, keypoints 3D và keypoints 2D. Trong trường hợp cần thiết, keypoints 2D có thể được sinh lại từ dữ liệu 3D thông qua phép chiếu camera.

Ở giai đoạn đầu tiên, hệ thống thực hiện đồng bộ frame giữa hai camera nhằm đảm bảo các frame tương ứng biểu diễn cùng một thời điểm chuyển động. Sau bước đồng bộ, chuỗi pose được đưa qua refinement\_pipeline để thực hiện các xử lý ban đầu như chuẩn hóa dữ liệu keypoints3D, tối ưu thông số nội tại của camera qua các giai đoạn tối ưu đồng thời loại bỏ các frame bất thường.

Tiếp theo, dữ liệu được đưa qua pose\_pipeline. Module này thực hiện các đánh giá với pose, shape là parameter trong mô hình SMPL sau đó cập nhật translation. Thực hiện duy trì tính nhất quán giữa các frame qua đánh giá offset. Sau bước xử lý pose, 2 tập dữ liệu tọa độ keypoints3D tương ứng của 2 view được sinh. Tiếp tục qua hệ thống sử dụng fusion\_pipeline để kết hợp thông tin từ hai góc nhìn nhằm giảm các sai lệch hình học và xử lý che khuất trong bài toán monocular pose.

Cuối cùng, dữ liệu sau fusion được đưa qua learnable\_pipeline. Pipeline này thực hiện tinh chỉnh các tham số SMPL theo hướng Learnable-SMPLify nhằm cải thiện sự phù hợp giữa pose reconstruction và dữ liệu quan sát. Kết quả sau cùng của hệ thống là chuỗi keypoints3D đã được tinh chỉnh để phục vụ cho quá trình đánh giá và trực quan hóa.

\subsection{Các baseline sử dụng để so sánh}

Do đồ án tập trung vào bài toán post-processing refinement vốn cần nhiều dữ liệu dự đoán từ model monocular, baseline phù hợp được sử dụng là kết quả thô từ mô hình WHAM trước khi áp dụng các bước hậu xử lý. Baseline này phản ánh chất lượng ban đầu của mô hình monocular pose estimation và được sử dụng để đánh giá mức độ cải thiện của pipeline refinement được đề xuất.

Ngoài baseline WHAM gốc, đồ án còn thực hiện so sánh giữa các giai đoạn trung gian trong pipeline nhằm phân tích ảnh hưởng của từng thành phần refinement. Cụ thể, kết quả sau refinement\_pipeline, pose\_pipeline, fusion\_pipeline và learnable\_pipeline được đánh giá độc lập nhằm xác định mức độ đóng góp của từng bước đối với chất lượng pose reconstruction cuối cùng.

Cách xây dựng baseline theo từng giai đoạn giúp đánh giá trực tiếp hiệu quả của các bước hậu xử lý mà không phụ thuộc hoàn toàn việc so sánh với mô hình monocular pose estimation. Điều này phù hợp với mục tiêu của đồ án vì trọng tâm nghiên cứu nằm ở quá trình refinement pose thay vì xây dựng hoặc huấn luyện một mô hình pose estimation mới.

\begin{table}[H]
\centering
\caption{Các cấu hình thực nghiệm so sánh}
\label{tab:ablation_design}
\begin{tabular}{|p{3.6cm}|p{4.8cm}|p{4.8cm}|}
\hline
\textbf{Cấu hình} & \textbf{Thành phần sử dụng} & \textbf{Mục tiêu đánh giá} \\
\hline
Baseline monocular & Output trực tiếp từ một view \cite{shin2024wham} & Chất lượng pose ban đầu từ mô hình monocular \\
\hline
Two-view aligned & Đồng bộ và căn chỉnh hai chuỗi pose \cite{gu2024multiview} & Ảnh hưởng của alignment đến sai khác liên góc nhìn \\
\hline
Consistency + Judgment & Alignment, consistency check và phân xử \cite{vu2026multiperspective} & Khả năng phát hiện và xử lý mâu thuẫn giữa hai view \\
\hline
Fusion + Temporal refinement & Fusion và ràng buộc thời gian \cite{physcap2020} & Mức giảm jitter và tăng ổn định chuyển động \\
\hline
Full pipeline & Toàn bộ các mô-đun đề xuất & Chất lượng sau phát hiện, phân xử và hiệu chỉnh đầy đủ \\
\hline
\end{tabular}
\end{table}

\subsection{Tiền xử lý dữ liệu}

Trước khi vào hệ thống, dữ liệu đầu vào được xử lý nhằm đồng nhất hệ tọa độ và giảm ảnh hưởng của nhiễu trong chuỗi pose thô.

Đầu tiên, hệ thống thực hiện đồng bộ frame giữa hai video nhằm đảm bảo tính tương ứng thời gian giữa các góc nhìn. Sau đó, các keypoints3D được chuẩn hóa về cùng hệ tọa độ để phục vụ cho quá trình fusion và đánh giá. Các frame có dữ liệu thiếu hoặc pose bất thường được loại bỏ nhằm hạn chế ảnh hưởng đến các bước tối ưu hóa phía sau.

Bên cạnh đó, các bước smoothing sơ bộ được áp dụng nhằm giảm rung giật trong chuỗi pose monocular ban đầu. Một số phép căn chỉnh rigid transformation cũng được sử dụng để đưa các pose về cùng không gian tham chiếu trước khi tính toán các độ đo đánh giá.

Ngoài dữ liệu pose, các keypoints 2D và keypoints 3D cũng được kiểm tra và đồng nhất cấu trúc skeleton nhằm đảm bảo tính tương thích giữa các thành phần trong pipeline xử lý.

\subsection{Thiết lập đánh giá}

Sau khi hoàn thành quá trình refinement, hệ thống tiến hành đánh giá chất lượng pose reconstruction bằng chỉ số MPJPE, PA-MPJPE và PCK. Các chỉ số này được tính trên keypoints 3D nhằm đo mức sai lệch giữa pose dự đoán và ground truth trong không gian ba chiều.

Để đánh giá ảnh hưởng của từng thành phần trong pipeline, các kết quả trung gian ở từng giai đoạn xử lý cũng được lưu lại và đánh giá độc lập. Việc so sánh giữa pose thô từ WHAM và keypoints3D sau refinement giúp phân tích khả năng cải thiện độ ổn định và tính hợp lý của chuyển động.

Bên cạnh đánh giá định lượng, hệ thống cũng thực hiện trực quan hóa kết quả dưới dạng skeleton 3D và chiếu ngược keypoints lên video 2D. Quá trình trực quan hóa được sử dụng để quan sát các lỗi thường gặp như rung giật chuyển động, sai lệch chiều sâu, foot sliding hoặc các pose không hợp lý về mặt giải phẫu.

Việc kết hợp giữa đánh giá định lượng và trực quan hóa giúp phản ánh đầy đủ hơn chất lượng của pipeline hệ thống, đặc biệt trong các trường hợp mà sai số hình học nhỏ nhưng chuyển động vẫn thiếu ổn định hoặc không tự nhiên.

\section{Đánh giá độ chính xác pose reconstruction}

\subsection{So sánh trước và sau hệ thống}

Bảng~\ref{tab:refinement_result} trình bày cấu trúc kết quả dùng để so sánh chất lượng pose trước và sau khi áp dụng pipeline refinement. Các ô giá trị hiện được để trống vì số liệu định lượng sẽ được cập nhật sau khi hoàn tất thực nghiệm.

\begin{table}[H]
\centering
\caption{Mẫu kết quả đánh giá pose reconstruction trước và sau refinement}
\label{tab:refinement_result}
\begin{tabular}{|p{3.7cm}|c|c|c|p{3.0cm}|}
\hline
\textbf{Chỉ số} & \textbf{Trước refinement} & \textbf{Sau refinement} & \textbf{Mức thay đổi} & \textbf{Kết luận} \\
\hline
MPJPE $\downarrow$ & -- & -- & -- & Sai số khớp trung bình \\
\hline
PA-MPJPE $\downarrow$ & -- & -- & -- & Sai số sau Procrustes alignment \\
\hline
PCK $\uparrow$ & -- & -- & -- & Tỷ lệ khớp đúng trong ngưỡng \\
\hline
Temporal jitter $\downarrow$ & -- & -- & -- & Độ ổn định theo thời gian \\
\hline
Bone-length variance $\downarrow$ & -- & -- & -- & Độ ổn định cấu trúc xương \\
\hline
Cross-view disagreement $\downarrow$ & -- & -- & -- & Mức nhất quán giữa hai view \\
\hline
Runtime/frame $\downarrow$ & -- & -- & -- & Khả năng triển khai thực tế \\
\hline
\end{tabular}
\end{table}

Các giá trị định lượng sẽ được cập nhật sau khi hoàn tất quá trình thực nghiệm trên tập dữ liệu đánh giá. Vì vậy, ở giai đoạn hiện tại, bảng chỉ đóng vai trò mô tả cấu trúc báo cáo kết quả và các chỉ số cần thu thập.

\subsection{Đánh giá trên các chuyển động phức tạp}

Đối với các chuyển động có tốc độ cao hoặc xuất hiện hiện tượng che khuất mạnh, pose estimation monocular thường xuất hiện các sai lệch về chiều sâu và rung giật theo thời gian.

Trong các trường hợp này, bước refinement được thiết kế để giảm các biến đổi bất thường giữa các frame liên tiếp và cải thiện tính liên tục của quỹ đạo chuyển động. Ngoài ra, việc sử dụng bổ sung thông tin đa góc nhìn hỗ trợ giảm ambiguity về chiều sâu đối với các khớp khó quan sát.

\section{Đánh giá ảnh hưởng của từng thành phần refinement}

\subsection{Ảnh hưởng của temporal smoothing}

Temporal smoothing được sử dụng nhằm giảm nhiễu và hạn chế rung giật giữa các frame liên tiếp.

Trong thực nghiệm, bước smoothing được kỳ vọng giúp cải thiện độ ổn định của quỹ đạo chuyển động, đặc biệt trong các chuỗi video có chuyển động nhanh hoặc thay đổi góc nhìn liên tục.

\subsection{Ảnh hưởng của multi-view refinement}

Multi-view refinement hỗ trợ hiệu chỉnh các sai lệch hình học phát sinh trong quá trình monocular pose estimation.

Việc khai thác bổ sung dữ liệu từ góc nhìn thứ hai giúp cải thiện độ chính xác của các khớp bị che khuất hoặc các chuyển động có ambiguity lớn về chiều sâu.

\subsection{Ảnh hưởng của các ràng buộc cơ sinh học}

Các ràng buộc động học và cơ sinh học giúp hạn chế các pose không hợp lý như biến dạng chiều dài xương hoặc xoắn khớp bất thường.

Trong quá trình thực nghiệm, các ràng buộc này được dùng để kiểm tra mức độ tăng tính hợp lý của chuyển động và mức giảm các lỗi xuất hiện trong các chuỗi chuyển động phức tạp.

\section{Đánh giá khả năng triển khai thực tế}

Hệ thống được đánh giá trên nhiều điều kiện chuyển động khác nhau bao gồm:

\begin{itemize}
    \item Chuyển động nhanh
    \item Hiện tượng che khuất
    \item Thay đổi góc nhìn lớn
    \item Chất lượng video không ổn định
\end{itemize}

Về mặt thiết kế, pipeline refinement hướng tới khả năng cải thiện tính ổn định của pose reconstruction trong nhiều tình huống thực tế mà không yêu cầu hệ thống đa camera phức tạp.

Ngoài ra, việc sử dụng các mô hình monocular pose estimation có sẵn giúp giảm yêu cầu phần cứng và chi phí triển khai của hệ thống so với một cấu hình đa camera phức tạp.

\section*{Kết thúc chương}

Chương này đã trình bày quy trình thực nghiệm và phương pháp đánh giá hệ thống refinement pose được đề xuất trong đồ án. Các chỉ số đánh giá, bộ dữ liệu sử dụng và các kịch bản thực nghiệm đã được mô tả nhằm phục vụ quá trình phân tích hiệu quả của hệ thống.

Bên cạnh đó, chương cũng trình bày khung đánh giá đối với độ chính xác pose reconstruction, độ ổn định chuyển động và ảnh hưởng của từng thành phần refinement trong pipeline đề xuất. Các nội dung này nhằm chuẩn bị cho phần báo cáo số liệu thực nghiệm, thay vì kết luận trước rằng phương pháp đã cải thiện chất lượng tái tạo chuyển động người 3D.

\end{document}
